% Generated draft: requires mathematical and source review.
\documentclass[11pt]{article}
\usepackage[margin=25mm]{geometry}
\usepackage{fontspec}
\setmainfont{Latin Modern Roman}
\usepackage{amsmath,amssymb,mathtools,unicode-math}
\setmathfont{Latin Modern Math}
\usepackage{enumitem,needspace,xurl,hyperref}
\setlength{\parindent}{0pt}
\setlength{\parskip}{5pt}
\setlength{\emergencystretch}{3em}
\newcommand{\N}{\mathbb{N}}
\newcommand{\Z}{\mathbb{Z}}
\newcommand{\Q}{\mathbb{Q}}
\newcommand{\R}{\mathbb{R}}
\newcommand{\C}{\mathbb{C}}
\newcommand{\F}{\mathbb{F}}
\newcommand{\A}{\mathbb{A}}
\newcommand{\PP}{\mathbb{P}}
\newcommand{\E}{\mathbb{E}}
\newcommand{\eps}{\varepsilon}
\newcommand{\dd}{\,\mathrm{d}}
\begin{document}
\section*{UnsolvedMath ID: 30000623}
\textbf{Source problem code:} catalogue entry
\subsection{Definitions and mathematical statement}
The source asks to extend its semi-infinite fusion construction from weights proportional to $\Lambda_0$ to every dominant weight, uniformly for the intended class of simple Lie algebras.

The notation $\Lambda_0$ suggests the usual affine framework: starting with a finite-dimensional complex simple Lie algebra $\mathfrak g$, form its associated affine Lie algebra $\widehat{\mathfrak g}$. Relative to affine simple coroots $h_i$, a dominant integral weight $\lambda$ satisfies $\lambda(h_i)\in\Z_{\geq0}$ for every $i$, and the zeroth fundamental weight satisfies $\Lambda_0(h_i)=\delta_{0i}$. The weights $m\Lambda_0$ are a distinguished subfamily of these weights. The affine type and normalization, including any degree-coordinate convention, must match the actual construction.

In that interpretation, the proposed extension has the quantifiers
\[
\text{for every intended simple }\mathfrak g
\text{ and every }\lambda\in P_+(\widehat{\mathfrak g}),
\quad\text{the source's semi-infinite fusion construction works at }\lambda.
\]
Here $P_+$ denotes dominant integral affine weights. This is a quantified description of the target, not a definition of ``works'': the extract supplies neither the finite fusion factors, filtration, connecting maps or semi-infinite limit, nor the module and structural properties that the construction is meant to realize. It also does not explicitly distinguish finite-dimensional from affine dominant weights or state the ground field, affine type and allowable levels. Thus the common affine interpretation is exposed as contextual rather than silently turned into a verified construction or a claimed module isomorphism.
\subsection{Short English statement}
Extend the intended semi-infinite fusion construction beyond multiples of the zeroth fundamental weight to every dominant weight for every intended simple Lie algebra. The construction, weight conventions and success criterion are missing from the extract.
\subsection{Sources}
\begin{itemize}
\item Supplied report: \emph{Infinite-Dimensional Lie Theory} (2007), Oberwolfach Reports, DOI 10.4171/owr/2006/55, pp. 3259--3320; proposed year 2006. \url{https://doi.org/10.4171/owr/2006/55}.
\item Local UnsolvedMath registry, permanent ID 30000623, catalogue entry; created 25 August 2026, updated 16 September 2026. Its corrected\_verified note describes editorial cleanup without mathematical additions. The supplied 21 August 2026 literature triage labels the extracted claim open and reports no located authoritative complete resolution. No additional external reference is supplied; no independent current-status assessment is asserted.
\end{itemize}
\end{document}
